\section{Neural-guided certificate generation}
\label{sec:neural-search}

The direct certificates were discovered by two separately trained models, one
for each combined quotient.  They share an architecture and training
procedure, but their state colorings, action sets, depth scales, parameters,
and checkpoints are distinct.  \Cref{tab:model-specifications} summarizes the
certificate-generation configurations.

\begin{table}[t]
\centering
\caption{Combined quotient and model specifications.}
\label{tab:model-specifications}
\small
\begin{tabular}{@{}lrr@{}}
\toprule
Property & Pair 3+4 & Pair 5+6\\
\midrule
Transition & $G_5\to G_7$ & $G_7\to G_9$\\
Combined quotient size & 14,234,011,545,600 & 1,664,617,163,489,280\\
State positions & 120 & 120\\
Position classes & 43 & 67\\
FTM actions & 28 & 20\\
Random-walk $K_{\max}$ & 22 & 26\\
Trainable parameters & 619,996 & 802,260\\
Certificate length limit & 21 & 25\\
\bottomrule
\end{tabular}
\end{table}

\subsection{Exact quotient encodings}

Each state is a coloring of the $120$ oriented sticker positions induced by
the target subgroup: positions movable inside the target are collapsed to one
class, while the remaining positions retain fixed classes.  The problem
builder uses Schreier--Sims to check both the subgroup index and equality
between the intended target subgroup and the pointwise stabilizer of the
retained classes~\cite{seress2003}.  A simultaneous conjugacy check aligns every generated
sticker action with the independent physical move convention.  Thus the model
sees the intended combined quotient rather than an accidental coarsening.
The committed problem manifests are the inputs to certificate replay; the
artifact recomputes their group orders and physical-action alignment without
loading a checkpoint.

The network first applies a position-specific class embedding.  Equivalently,
for coloring $x=(c_1,\ldots,c_{120})$ it computes
\(
 e(x)=b+\sum_{p=1}^{120}W_{p,c_p}
\)
without materializing the full position--class one-hot vector.  Batch
normalization and the SiLU activation map this $64$-dimensional representation
to width $256$.  Each of two residual blocks has the form
\[
 h\leftarrow\operatorname{SiLU}\!\left(
 h+\operatorname{BN}_2\!\left(L_2(
 \operatorname{Dropout}(\operatorname{SiLU}(\operatorname{BN}_1(L_1(h)))))
 \right)\right).
\]
Dropout was zero in both published runs.  A final linear layer emits one
real-valued score per legal action; lower scores are searched first.

\subsection{Sparse trajectory supervision}
\label{sec:sparse-supervision}

Training examples are generated online by nontrivial random walks starting at
the target coordinate.  Consecutive moves on the same face are excluded
because their product is either one FTM move or the identity; quotient
self-loops are also rejected.  For a randomly selected trajectory pivot $x_t$,
let $a^-$ be the inverse of the preceding action and $a^+$ the sampled next
action.  Only these two outputs are supervised:
\begin{equation}
  \mathcal{L}(x_t)
  =\frac12\left(Q_\theta(x_t,a^-)-(t-1)\right)^2
  +\frac12\left(Q_\theta(x_t,a^+)-(t+1)\right)^2.
  \label{eq:sparse-loss}
\end{equation}
The diagnostic pair accuracy is the fraction for which
$Q_\theta(x_t,a^-)<Q_\theta(x_t,a^+)$.  The trajectory index $t$ is not asserted
to equal the shortest distance of $x_t$.  Accordingly, $Q_\theta$ is best
understood as a learned action-ranking score trained from noisy trajectory
depths, not as an exact distance table or a certified Q-function.

\begin{table}[t]
\centering
\caption{Training configuration used for both models, except for the
pair-specific random-walk limit.}
\label{tab:training-configuration}
\small
\begin{tabular}{@{}ll@{\qquad}ll@{}}
\toprule
Optimizer & AdamW & Learning rate & $10^{-4}$\\
Weight decay & $0.003$ & Gradient clip & $1$\\
Batch size & $1{,}024$ & Steps per epoch & $256$\\
Walk range & $2\ldots K_{\max}$ & Validation size & $16{,}384$\\
Seed & $42$ & Inference/training type & bfloat16\\
\bottomrule
\end{tabular}
\end{table}

For every sample, the walk length is uniform on the displayed integer range
and the pivot is uniform among its internal vertices.  Eligible nontrivial
moves are sampled uniformly after the same-face exclusion.  The pair-3+4
certificate run used epochs 1024 and 1120 of one training run; pair 5+6 used
epochs 864 and 3744.  These checkpoint identities record the run that found
the certificates, not a claim that their epochs are generally optimal.

\subsection{One beam per root, evaluated in batches}
\label{sec:batched-beam}

At each depth, every root maintains its own beam, visited set, last-face
context, and paths; only neural inference is shared across roots.  Scores are
the current one-step values $Q_\theta(x,a)$, not cumulative path scores.  At a
beam of width $B$ with $A$ actions, the implementation first obtains the best
\[
  C=\min\{BA,\max(A,16B)\}
\]
parent--action pairs before CPU deduplication.  It initially scans
\(\min\{C,\max(A,2B)\}\) candidates and scans the rest only if duplicates
prevent the next beam from filling.  Same-face continuations are masked.  The
goal test occurs immediately after applying an action; self-loops and then
previously visited keys are discarded.  A key is the pair
\((\text{coloring},\text{last face})\), because those two copies of one
coloring have different legal continuations.  The first $B$ surviving keys in
score order form the next beam.

\begin{lstlisting}[caption={One root of the batched beam search.  Model calls
for all active roots are fused in the implementation.},label={lst:inner-beam}]
beam := {(root, last_face = none, word = empty)}
visited := {(root, none)}
for depth := 1 .. length_limit:
    P := legal (parent, action) pairs from beam
    rank P by one-step score Q(parent, action), ascending
    next := empty
    for pair in bounded top-candidate window P:
        child := apply(pair)
        if child is target: return pair.word
        key := (child, face(pair.action))
        if child is a self-loop or key in visited: continue
        add key to visited and pair to next
        if |next| = beam_width: break
    if next is empty: return failure
    beam := next
return failure
\end{lstlisting}

Exact floating-point ties inherit whatever ordering the PyTorch device
produces; the search does not require that order to be reproducible.  Any returned word is replayed
in the combined quotient and on the corresponding physical state before it can
enter the certificate database.

\begin{lstlisting}[caption={Survivor beam cascade used for certificate discovery.},
                   label={lst:cascade}]
U := ordered set of unresolved frontier IDs
for (checkpoint, beam_width) in schedule:
    for bounded chunk C of U:
        words := batched_per_root_beams(C, checkpoint,
                                       beam_width, length_limit)
        for (id, word) in words:
            if replay_quotient(id, word) and replay_physical(id, word):
                persist_certificate(id, word, provenance)
    U := U minus all persisted certificate IDs
    stop if U is empty
require U is empty
\end{lstlisting}

The database is resumable: existing IDs are skipped, records are tied to
checksums of the physical state and proof problem, and a later stage may replace
a record only with a shorter word or the same length at a smaller beam width.
The production runner uses outer chunks of $10{,}000$ frontier IDs and root
minibatches that decrease from 128 at narrow widths to 16 at width 4096.  The
older pair-5+6 width-64 pass used an earlier low-batch implementation, which is
why its throughput is not compared directly with later tiers.  Beam width,
checkpoint identity, and batching record discovery provenance only.

\subsection{Published cascade schedules}

The survivor plot in \cref{fig:survivor-curves} shows the schedules; full
counts, timings, and checkpoint splits appear in \cref{app:cascade}.  Roughly
60\% of each direct-certificate set was found at its first beam tier, while
the maximum width was applied to only one pair-3+4 survivor and $28$
pair-5+6 survivors.  These observations do not isolate the causal effect of
beam width.  Pair 3+4 combines
epoch-1024 and epoch-1120 checkpoints at width 32; pair 5+6 switches from
epoch 864 at width 64 to epoch 3744 for subsequent tiers.  Moreover, every
later tier sees a survivor distribution selected by all earlier tiers.  We
therefore present the cascade as the successful certificate-generation
portfolio actually used, not as evidence that one beam width dominates
another.
